\chapter{Descomposición en valores singulares (SVD)}\label{cap:SVD}

Las matrices simétricas son cómodas: tienen autovalores reales y una base
ortonormal de autovectores, y eso es todo el Capítulo~\ref{cap:schur}. Las
matrices cuadradas generales ya no: los autovalores pueden ser complejos y
los autovectores pueden no ser ortogonales, o no alcanzar para formar una
base. Y si $A$ no es cuadrada, la ecuación $A\bx=\lambda\bx$ ni siquiera tiene
sentido, porque el miembro izquierdo vive en $\R^m$ y el derecho en $\R^n$.

Hace falta una idea que funcione para \emph{cualquier} matriz, y la
descomposición en valores singulares la provee de una manera que resulta
perfecta. El precio a pagar es abandonar la idea de una sola base: en lugar de
una familia de autovectores necesitaremos \textbf{dos} familias ortonormales,
una en $\R^n$ y otra en $\R^m$, relacionadas no por $A\bx=\lambda\bx$ sino por
\[
A\bv_k=\sigma_k\bu_k .
\]
La matriz manda cada $\bv_k$ a un múltiplo de $\bu_k$: no conserva la
dirección --- no puede, si cambia de espacio --- pero manda una base
ortonormal a otra base ortonormal, estirando cada eje por su propio factor.

El capítulo se organiza así. En la Sección~\ref{sec:sigma1} definimos los
valores singulares \emph{antes} de tener la SVD, a partir de la matriz
simétrica $A^TA$, y probamos que el primero de ellos es la norma que mide
cuánto estira $A$. Recién después (Secciones~\ref{sec:SVD}
y~\ref{sec:prueba-SVD}) enunciamos y probamos la descomposición. El resto del
capítulo es leer lo que la SVD dice: bases ortonormales de los cuatro
subespacios fundamentales, la imagen de la esfera unitaria, y la
descomposición polar.

Dos resultados centrales quedan para después. El \textbf{teorema de
Eckart-Young} --- que truncar la SVD da la mejor aproximación de rango bajo
--- se prueba en el capítulo de PCA, que es donde se usa. Y la
\textbf{pseudoinversa}, que la SVD permite definir para cualquier matriz, se
trata en el capítulo final junto con los cuadrados mínimos.

% ===================================================================
\section{Valores singulares y la norma espectral}
\label{sec:sigma1}

Sea $A\in\R^{m\times n}$, de rango $r$. Aunque $A$ no sea cuadrada, la matriz
$A^TA\in\R^{n\times n}$ siempre es simétrica y semidefinida positiva:
\[
(A^TA)^T=A^TA,
\qquad
\bx^T(A^TA)\bx=\|A\bx\|^2\ge0 .
\]
Por el teorema espectral (Teorema~\ref{teo:espectral-real}) tiene autovalores
reales, y por la Entrega~\ref{ent:semidefinidas} todos son $\ge0$. Además,
por la Observación~\ref{obs:ATA} sabemos que $N(A^TA)=N(A)$, de modo que
$\rg(A^TA)=\rg(A)=r$: exactamente $r$ de esos autovalores son positivos.

\begin{definition}[Valores singulares]\label{def:valores-singulares}
Sean $\lambda_1\ge\lambda_2\ge\cdots\ge\lambda_n\ge0$ los autovalores de
$A^TA$, y sean $\bv_1,\dots,\bv_n$ autovectores ortonormales
correspondientes. Los \textbf{valores singulares} de $A$ son
\[
\sigma_k=\sqrt{\lambda_k},\qquad k=1,\dots,n,
\]
de modo que $\sigma_1\ge\sigma_2\ge\cdots\ge\sigma_r>0$ y
$\sigma_{r+1}=\cdots=\sigma_n=0$. Los $\bv_k$ se llaman \textbf{vectores
singulares a derecha}.
\end{definition}

Conviene subrayar que esto es una \emph{definición}, no un teorema: no
necesitamos la SVD para hablar de los $\sigma_k$. Lo que sigue tampoco.

\begin{theorem}[El primer valor singular]\label{teo:sigma1}
Sea $A\in\R^{m\times n}$ no nula. Entonces
\[
\max_{\|\bx\|=1}\|A\bx\| = \sigma_1 ,
\qquad\text{equivalentemente}\qquad
\max_{\bx\neq0}\frac{\|A\bx\|}{\|\bx\|}=\sigma_1 ,
\]
y el máximo se alcanza exactamente en los autovectores unitarios de $A^TA$
asociados a $\lambda_1$; en particular en $\bx=\bv_1$.
\end{theorem}

\begin{proof}
Conviene normalizar primero y elevar al cuadrado después. Si $\|\bx\|=1$,
\[
\|A\bx\|^2=(A\bx)^T(A\bx)=\bx^T(A^TA)\bx ,
\]
de modo que maximizar $\|A\bx\|$ sobre los vectores unitarios es maximizar la
\emph{energía} de la matriz simétrica $S=A^TA$ sobre la esfera. Por el
Teorema~\ref{teo:rayleigh}, ese máximo vale $\lambda_1$ ---el mayor autovalor
de $A^TA$--- y se alcanza exactamente en los autovectores unitarios asociados
a $\lambda_1$. Tomando raíz cuadrada,
$\max_{\|\bx\|=1}\|A\bx\|=\sqrt{\lambda_1}=\sigma_1$.

La versión con cociente es la misma, por el
Corolario~\ref{coro:rayleigh-cociente}: para $\bx\neq0$ cualquiera,
$\|A\bx\|/\|\bx\|=\|A\bu\|$ con $\bu=\bx/\|\bx\|$ unitario.
\end{proof}

El teorema dice que el cociente $\|A\bx\|/\|\bx\|$ --- cuánto estira $A$ al
vector $\bx$ --- está acotado, y que $\sigma_1$ es el estiramiento máximo. Eso
es exactamente lo que hace falta para medir una matriz.

\begin{definition}[Norma espectral]\label{def:norma-espectral}
La \textbf{norma espectral} de $A\in\R^{m\times n}$ es
\[
\boxed{\ \|A\|_2=\max_{\bx\neq0}\frac{\|A\bx\|}{\|\bx\|}
=\max_{\|\bx\|=1}\|A\bx\|=\sigma_1(A)\ }
\]
\end{definition}

\begin{proposition}\label{prop:espectral-norma}
$\|\cdot\|_2$ es una norma en $\R^{m\times n}$, y para todo $\bx\in\R^n$ vale
\[
\|A\bx\|\le\|A\|_2\,\|\bx\| .
\]
\end{proposition}

\begin{proof}
La desigualdad es la definición del máximo. Para los axiomas de norma
(Sección~\ref{sec:normas}):

\emph{(n2)} Claramente $\|A\|_2\ge0$. Si $\|A\|_2=0$ entonces $A\bx=0$ para
todo $\bx$, en particular para los $\bi{e}_j$, así que todas las columnas de $A$
son nulas y $A=0$.

\emph{(n1)} $\|cA\|_2=\max_{\|\bx\|=1}\|cA\bx\|=|c|\max_{\|\bx\|=1}\|A\bx\|
=|c|\,\|A\|_2$.

\emph{(n3)} Para $\|\bx\|=1$, usando la desigualdad triangular de la norma de
vectores,
\[
\|(A+B)\bx\|\le\|A\bx\|+\|B\bx\|\le\|A\|_2+\|B\|_2 ,
\]
y tomando el máximo sobre esos $\bx$ queda $\|A+B\|_2\le\|A\|_2+\|B\|_2$.
\end{proof}

\begin{remark}[Dos normas de matrices, dos ideas distintas]
Ya tenemos dos maneras de medir una matriz. La de Frobenius
(Sección~\ref{sec:frobenius}) se olvida de que $A$ es una matriz y la mide
como al vector largo de sus $mn$ entradas; viene de un producto interno. La
espectral la mide \emph{como transformación}: cuánto estira en el peor caso.
Son distintas y las dos se usan. En el capítulo de PCA se verá cómo se
relacionan con los $\sigma_k$ y entre sí; en el
capítulo de normas de matrices estudiaremos la familia general de normas
definidas por un máximo como el de arriba, y qué se gana al pedirle a una
norma de matrices que se lleve bien con el producto.
\end{remark}

\begin{corollary}[Los autovalores no superan a $\sigma_1$]\label{coro:lambda-sigma}
Si $A\in\R^{n\times n}$ y $\lambda$ es un autovalor real de $A$, entonces
$|\lambda|\le\sigma_1$.
\end{corollary}

\begin{proof}
Si $A\bx=\lambda\bx$ con $\bx\neq0$, entonces
$|\lambda|\,\|\bx\|=\|A\bx\|\le\sigma_1\|\bx\|$.
\end{proof}

\begin{example}\label{ej:rango1-sigma}
Para una matriz de rango uno $A=\bx\by^T$ (con $\bx,\by$ no nulos), la
desigualdad anterior es Cauchy-Schwarz disfrazada. En efecto, $A^TA
=\by\bx^T\bx\by^T=\|\bx\|^2\,\by\by^T$ tiene autovalor no nulo
$\|\bx\|^2\|\by\|^2$, de modo que
\[
\sigma_1=\|\bx\|\,\|\by\| ,
\qquad
\bv_1=\frac{\by}{\|\by\|},
\qquad
\bu_1=\frac{\bx}{\|\bx\|} .
\]
Por otro lado, el único autovalor no nulo de $A$ es $\lambda=\by^T\bx$, con
autovector $\bx$, porque $(\bx\by^T)\bx=\bx(\by^T\bx)$. El
Corolario~\ref{coro:lambda-sigma} dice entonces
$|\by^T\bx|\le\|\bx\|\,\|\by\|$.
\end{example}

Necesitaremos también la versión del Teorema~\ref{teo:sigma1} para los
valores singulares siguientes, que se obtiene restringiendo el máximo.

\begin{corollary}[Los $\sigma_k$ siguientes]\label{coro:sigma-k}
Con las notaciones de la Definición~\ref{def:valores-singulares}, para cada
$k=1,\dots,n$,
\[
\max\bigl\{\|A\bx\| : \|\bx\|=1,\ \bx\perp\bv_1,\dots,\bv_{k-1}\bigr\}=\sigma_k ,
\]
y el máximo se alcanza en $\bx=\bv_k$.
\end{corollary}

\begin{proof}
Escribamos $\bx=\sum_i c_i\bv_i$ en la base ortonormal de autovectores. La
condición $\bx\perp\bv_1,\dots,\bv_{k-1}$ dice $c_1=\cdots=c_{k-1}=0$, y
$\|\bx\|=1$ dice $\sum_i c_i^2=1$. Entonces
\[
\|A\bx\|^2=\bx^T(A^TA)\bx=\sum_{i\ge k}\lambda_i c_i^2
\le\lambda_k\sum_{i\ge k}c_i^2=\lambda_k ,
\]
con igualdad cuando $c_k=1$ y el resto cero, es decir en $\bx=\bv_k$.
\end{proof}

Este corolario es la manera de entender los vectores singulares \textbf{uno
por uno} en lugar de todos juntos: $\bv_1$ es la dirección que más se estira;
$\bv_2$ es la que más se estira entre las perpendiculares a $\bv_1$; y así
sucesivamente. Es también el punto de partida del capítulo de PCA, donde este
argumento se generaliza a $d$ direcciones a la vez.

% ===================================================================
\section{La descomposición en valores singulares}
\label{sec:SVD}

Ya tenemos la mitad de los ingredientes: los $\sigma_k$ y los $\bv_k$. Los
vectores singulares a izquierda salen ahora de la relación que queremos que
valga.

\begin{definition}[Vectores singulares a izquierda]\label{def:u}
Para $k=1,\dots,r$ definimos
\[
\bu_k=\frac{A\bv_k}{\sigma_k}\in\R^m .
\]
\end{definition}

Con esa definición, $A\bv_k=\sigma_k\bu_k$ vale por construcción para
$k\le r$; y para $k>r$ es $\|A\bv_k\|^2=\lambda_k=0$, o sea $A\bv_k=\bi 0$.
Es decir:
\begin{equation}\label{eq:Av-sigma-u}
\boxed{\ A\bv_1=\sigma_1\bu_1,\ \ \dots,\ \ A\bv_r=\sigma_r\bu_r\ }
\qquad
\boxed{\ A\bv_{r+1}=\bi 0,\ \ \dots,\ \ A\bv_n=\bi 0\ }
\end{equation}
Los últimos $n-r$ vectores $\bv$ están en el espacio nulo de $A$.

Falta ver que los $\bu_k$ son ortonormales --- lo haremos en la sección
siguiente --- y completar la lista hasta $\bu_m$. Aceptándolo por un momento,
pasemos~\eqref{eq:Av-sigma-u} a forma matricial. Son $n$ ecuaciones, una por
cada $\bv_k$, y lo que queremos es escribirlas como \emph{una sola} igualdad
entre matrices. La herramienta es la fórmula de las columnas de un producto
del Capítulo~\ref{cap:algebra-lineal}.

Pongamos los $\bv$ como columnas de $V\in\R^{n\times n}$ y los $\bu$ como
columnas de $U\in\R^{m\times m}$, ambas ortogonales, y llamemos
$\Sigma\in\R^{m\times n}$ a la matriz con $\sigma_1,\dots,\sigma_r$ en las
primeras posiciones diagonales y ceros en todo lo demás. Ahora comparemos las
columnas de $AV$ con las de $U\Sigma$, una por una.

\medskip
\noindent\textbf{La columna $j$ de $AV$.} Por~\eqref{eq:cols},
\[
\col_j(AV)=A\,\col_j(V)=A\bv_j .
\]

\noindent\textbf{La columna $j$ de $U\Sigma$.} Otra vez por~\eqref{eq:cols},
$\col_j(U\Sigma)=U\,\col_j(\Sigma)$. Y las columnas de $\Sigma$ son fáciles:
$\col_j(\Sigma)=\sigma_j\bi{e}_j$ si $j\le r$, y $\col_j(\Sigma)=\bi 0$ si
$j>r$. Usando ahora la fórmula del producto matriz por vector
\eqref{eq:mat-vec-col}, que dice que $U\bx$ es la combinación de las columnas
de $U$ con coeficientes las entradas de $\bx$,
\[
\col_j(U\Sigma)=U(\sigma_j\bi{e}_j)=\sigma_j\,\col_j(U)=\sigma_j\bu_j
\quad (j\le r),
\qquad
\col_j(U\Sigma)=U\bi 0=\bi 0\quad (j>r).
\]

\medskip
Comparando las dos listas con~\eqref{eq:Av-sigma-u}: para $j\le r$ las dos
columnas valen $\sigma_j\bu_j$, y para $j>r$ las dos valen $\bi 0$. Dos
matrices con las mismas columnas son la misma matriz, así que
\begin{equation}\label{eq:AV-USigma}
\boxed{\ AV=U\Sigma\ }
\end{equation}
Las $n$ ecuaciones de~\eqref{eq:Av-sigma-u} quedaron empaquetadas en una. La
Figura~\ref{fig:esquema-SVD} muestra los tamaños y qué hay en cada factor.

\begin{figure}[htbp]
\centering
\begin{tikzpicture}[scale=0.52, every node/.style={font=\footnotesize}]
  \def\m{6} \def\n{4} \def\r{2.4}
  % ---------------- A (m x n) ----------------
  \draw[fill=bleu!8,draw=bleu!70!black,thick] (0,0) rectangle (\n,\m);
  \node at (\n/2,\m/2) {$A$};
  \node[below] at (\n/2,-0.2) {$m\times n$};
  \node at (4.75,\m/2) {$\cdot$};
  % ---------------- V (n x n) ----------------
  \begin{scope}[xshift=5.5cm]
    \draw[fill=gray!8,draw=black!65,thick] (0,\m-\n) rectangle (\n,\m);
    \draw[draw=black!40] (\r,\m-\n)--(\r,\m);
    \node at (\r/2,\m-\n/2) {$V_r$};
    \node at ({(\r+\n)/2},\m-\n/2) {$V_0$};
    \node[below,align=center] at (\n/2,\m-\n-0.2) {$n\times n$\\ ortogonal};
  \end{scope}
  \node[scale=1.5] at (10.6,\m/2) {$=$};
  % ---------------- U (m x m) ----------------
  \begin{scope}[xshift=11.7cm]
    \draw[fill=gray!8,draw=black!65,thick] (0,0) rectangle (\m,\m);
    \draw[draw=black!40] (\r,0)--(\r,\m);
    \node at (\r/2,\m/2) {$U_r$};
    \node at ({(\r+\m)/2},\m/2) {$U_0$};
    \node[below,align=center] at (\m/2,-0.2) {$m\times m$\\ ortogonal};
  \end{scope}
  \node at (18.4,\m/2) {$\cdot$};
  % ---------------- Sigma (m x n) ----------------
  \begin{scope}[xshift=19.2cm]
    \draw[fill=gray!14,draw=black!65,thick] (0,0) rectangle (\n,\m);
    \fill[orange!20] (0,\m-\r) rectangle (\r,\m);
    \draw[orange!70!black,thick] (0,\m-\r) rectangle (\r,\m);
    \draw[orange!70!black] (0.15,\m-0.15)--({\r-0.15},{\m-\r+0.15});
    \node[orange!45!black,scale=0.8] at (0.92,\m-0.62) {$\sigma_1$};
    \node[orange!45!black,scale=0.8] at (1.92,\m-1.68) {$\sigma_r$};
    \node[black!50] at ({(\r+\n)/2},\m-\r/2) {$0$};
    \node[black!50] at (\n/2,{(\m-\r)/2}) {$0$};
    \node[below] at (\n/2,-0.2) {$m\times n$};
  \end{scope}
\end{tikzpicture}
\caption{La igualdad $AV=U\Sigma$ en tamaños. Las columnas de $V$ y de $U$
vienen partidas en dos grupos: $V_r$ son las primeras $r$ columnas de $V$
---una base ortonormal del espacio fila $C(A^T)$--- y $V_0$ las restantes,
que son una base de $N(A)$; $U_r$ es una base de $C(A)$ y $U_0$ una de
$N(A^T)$. La matriz $\Sigma$ tiene los $\sigma_k$ en el bloque $r\times r$ de
arriba a la izquierda y ceros en todo lo demás: es la que ``recorta'', y la
que explica por qué de $U$ sólo llegan a usarse las primeras $r$ columnas.}
\label{fig:esquema-SVD}
\end{figure}

Y como $V$ es ortogonal, $V^{-1}=V^T$, y multiplicando~\eqref{eq:AV-USigma} a
derecha por $V^T$:

\begin{theorem}[Descomposición en valores singulares]\label{teo:SVD}
Toda matriz $A\in\R^{m\times n}$ se puede escribir como
\[
\boxed{\ A=U\,\Sigma\,V^T\ }
\]
con $U\in\R^{m\times m}$ y $V\in\R^{n\times n}$ ortogonales y
$\Sigma\in\R^{m\times n}$ ``diagonal'' con entradas
$\sigma_1\ge\cdots\ge\sigma_r>0$ y el resto ceros.
\end{theorem}

Para ver qué hay adentro de ese producto volvemos al Capítulo~\ref{cap:algebra-lineal},
esta vez a la fórmula~\eqref{eq:suma-rango1}, que escribe un producto de
matrices como suma de matrices de rango uno. Agrupemos $A=(U\Sigma)\,V^T$ y
apliquémosla al par $U\Sigma$ (de $m\times n$) y $V^T$ (de $n\times n$):
\[
A=\sum_{p=1}^{n}\col_p(U\Sigma)\ \fila_p(V^T) .
\]
Los dos factores de cada sumando ya los calculamos: $\col_p(U\Sigma)$ vale
$\sigma_p\bu_p$ si $p\le r$ y $\bi 0$ si $p>r$, y $\fila_p(V^T)=\bv_p^T$. Los
$n-r$ últimos sumandos son nulos, y queda
\begin{equation}\label{eq:sumandos}
\boxed{\ A=U\Sigma V^T
=\sigma_1\bu_1\bv_1^T+\sigma_2\bu_2\bv_2^T+\cdots+\sigma_r\bu_r\bv_r^T\ }
\end{equation}

\begin{remark}[Por qué importa esta forma]
Como $A=LU$, como $A=QR$ y como $S=Q\Lambda Q^T$, la SVD separa la matriz en
sumandos de rango uno. Lo que la distingue es que \emph{los sumandos vienen
ordenados por importancia}: $\sigma_1\bu_1\bv_1^T$ es la matriz de rango uno
más cercana a $A$, y más en general, la suma de los primeros $d$ sumandos es
la mejor aproximación de rango $d$ que existe. Ése es el teorema de
Eckart-Young, y es el motivo por el que la SVD es la factorización central en
ciencia de datos. Lo enunciaremos con precisión en la
Sección~\ref{sec:lo-que-viene} y se prueba en el capítulo de PCA.
\end{remark}

\subsection{La forma reducida}

En~\eqref{eq:sumandos} se ve que los últimos $n-r$ sumandos no aportan nada:
el corazón de la SVD está en los primeros $r$ vectores. Eso sugiere quedarse
sólo con ellos. Llamemos
\[
U_r=[\,\bu_1\ \cdots\ \bu_r\,]\in\R^{m\times r},
\qquad
V_r=[\,\bv_1\ \cdots\ \bv_r\,]\in\R^{n\times r},
\qquad
\Sigma_r=\diag(\sigma_1,\dots,\sigma_r) .
\]

Lo interesante es que para llegar a la forma reducida no hace falta ninguna
cuenta nueva: alcanza con leer~\eqref{eq:suma-rango1} \emph{al revés}. Esa
fórmula dice que una suma de productos externos es un producto de matrices,
con los vectores de la izquierda como columnas del primer factor y los de la
derecha como filas del segundo. Y en~\eqref{eq:sumandos} tenemos justamente
una suma de $r$ productos externos:
\[
A=\sum_{p=1}^{r}(\sigma_p\bu_p)\,\bv_p^T
=\sum_{p=1}^{r}\col_p(U_r\Sigma_r)\ \fila_p(V_r^T) ,
\]
donde usamos que $\col_p(U_r\Sigma_r)=\sigma_p\bu_p$ ---la misma cuenta de
antes, ahora con $\Sigma_r$ cuadrada--- y que $\fila_p(V_r^T)=\bv_p^T$.
Aplicando~\eqref{eq:suma-rango1} de derecha a izquierda:
\[
\boxed{\ A=U_r\,\Sigma_r\,V_r^T\ }\qquad\text{(\textbf{SVD reducida}).}
\]
La Figura~\ref{fig:esquema-SVD-reducida} la muestra al lado de la completa.

\begin{figure}[htbp]
\centering
\begin{tikzpicture}[scale=0.52, every node/.style={font=\footnotesize}]
  \def\m{6} \def\n{4} \def\r{2.4}
  % --------- A ---------
  \draw[fill=bleu!8,draw=bleu!70!black,thick] (0,0) rectangle (\n,\m);
  \node at (\n/2,\m/2) {$A$};
  \node[below] at (\n/2,-0.15) {$m\times n$};
  \node[scale=1.6] at (5.1,\m/2) {$=$};
  % --------- U_r (m x r) ---------
  \begin{scope}[xshift=6.3cm]
    \draw[fill=gray!6,draw=black!65,thick] (0,0) rectangle (\r,\m);
    \node at (\r/2,\m/2) {$U_r$};
    \node[below] at (\r/2,-0.15) {$m\times r$};
    \node[above,align=center] at (\r/2,\m+0.2) {base de $C(A)$};
  \end{scope}
  \node at (9.3,\m/2) {$\cdot$};
  % --------- Sigma_r (r x r) ---------
  \begin{scope}[xshift=10.1cm]
    \fill[orange!22,draw=orange!70!black,thick] (0,\m-\r) rectangle (\r,\m);
    \draw[orange!70!black] (0.15,\m-0.15)--({\r-0.15},{\m-\r+0.15});
    \node[orange!45!black,scale=0.8] at (0.92,\m-0.62) {$\sigma_1$};
    \node[orange!45!black,scale=0.8] at (1.92,\m-1.68) {$\sigma_r$};
    \node[below,align=center] at (\r/2,\m-\r-0.2) {$r\times r$\\ diagonal};
  \end{scope}
  \node at (13.1,\m/2) {$\cdot$};
  % --------- V_r^T (r x n) ---------
  \begin{scope}[xshift=13.9cm]
    \draw[fill=gray!6,draw=black!65,thick] (0,\m-\r) rectangle (\n,\m);
    \node at (\n/2,\m-\r/2) {$V_r^T$};
    \node[below] at (\n/2,\m-\r-0.15) {$r\times n$};
    \node[above,align=center] at (\n/2,\m+0.2) {base de $C(A^T)$, por filas};
  \end{scope}
\end{tikzpicture}
\caption{La SVD reducida $A=U_r\Sigma_r V_r^T$. Comparada con la
Figura~\ref{fig:esquema-SVD}, desaparecieron las columnas de $U$ y de $V$ que
$\Sigma$ multiplicaba por cero, y $\Sigma_r$ quedó cuadrada e invertible.
Guardar $A$ cuesta $mn$ números; guardar los tres factores de la derecha,
$r(m+n+1)$. Cuando $r$ es chico frente a $m$ y $n$, la diferencia es enorme:
es la idea que explota el capítulo de PCA.}
\label{fig:esquema-SVD-reducida}
\end{figure}

También se puede llegar por el otro lado, que es el camino del
Ejercicio~\ref{ej:reducida}: las ecuaciones~\eqref{eq:Av-sigma-u} que no son
triviales son las $r$ columnas de
\begin{equation}\label{eq:SVD-reducida}
A\,\underbrace{[\,\bv_1\ \cdots\ \bv_r\,]}_{\text{base del espacio fila}}
=\underbrace{[\,\bu_1\ \cdots\ \bu_r\,]}_{\text{base del espacio columna}}
\Sigma_r ,
\qquad\text{es decir}\qquad
AV_r=U_r\Sigma_r ,
\end{equation}
y de ahí se multiplica por $V_r^T$. Pero atención: acá $V_r$ \emph{no} es
cuadrada, así que $V_r^T$ no es su inversa y el paso necesita un argumento
---es justamente lo que pide ese ejercicio---, mientras que por el camino de
los productos externos la igualdad sale sin invertir nada.

Sigue valiendo $V_r^TV_r=I_r$ y $U_r^TU_r=I_r$, porque las columnas son
ortonormales. Pero cuando $V_r$ y $U_r$ no son cuadradas ya no hay inversa de
los dos lados: en general $V_rV_r^T\neq I$ y $U_rU_r^T\neq I$.

\begin{example}
Con $V_r=\begin{bsmallmatrix}1/3\\ 2/3\\ 2/3\end{bsmallmatrix}$ es
$V_r^TV_r=[1]$, pero
\[
V_rV_r^T=\frac19\begin{bmatrix}1&2&2\\ 2&4&4\\ 2&4&4\end{bmatrix}
\]
tiene rango $1$. (Lo que sí es cierto, y lo veremos en la
Sección~\ref{sec:cuatro-SVD}, es que $V_rV_r^T$ es la proyección ortogonal
sobre el espacio fila de $A$.)
\end{example}

\subsection{Un ejemplo completo}

\begin{example}\label{ej:svd-completo}
Calculemos $U$, $\Sigma$ y $V$ para
$A=\begin{bsmallmatrix}3&0\\ 4&5\end{bsmallmatrix}$.

El rango es $2$, así que habrá dos valores singulares positivos. Empezamos por
las dos matrices simétricas:
\[
A^TA=\begin{bmatrix}25&20\\ 20&25\end{bmatrix},
\qquad
AA^T=\begin{bmatrix}9&12\\ 12&41\end{bmatrix} .
\]
Tienen la misma traza ($50$) y los mismos autovalores, $\lambda_1=45$ y
$\lambda_2=5$; luego $\sigma_1=\sqrt{45}=3\sqrt5$ y $\sigma_2=\sqrt5$.
Observemos de paso que $\sigma_1\sigma_2=3\sqrt5\cdot\sqrt5=15=\det A$, y que
$\sigma_1>5=\lambda_{\max}(A)$ mientras que $\sigma_2<3=\lambda_{\min}(A)$:
los valores singulares \emph{no} son los autovalores.

Los autovectores de $A^TA$ son inmediatos:
\[
\begin{bmatrix}25&20\\ 20&25\end{bmatrix}\begin{pmatrix}1\\ 1\end{pmatrix}
=45\begin{pmatrix}1\\ 1\end{pmatrix},
\qquad
\begin{bmatrix}25&20\\ 20&25\end{bmatrix}\begin{pmatrix}-1\\ 1\end{pmatrix}
=5\begin{pmatrix}-1\\ 1\end{pmatrix},
\]
y normalizando,
$\bv_1=\frac1{\sqrt2}(1,1)^T$, $\bv_2=\frac1{\sqrt2}(-1,1)^T$.
Los $\bu$ salen de la Definición~\ref{def:u}:
\[
A\bv_1=\frac{3}{\sqrt2}\begin{pmatrix}1\\ 3\end{pmatrix}
=\sqrt{45}\cdot\frac1{\sqrt{10}}\begin{pmatrix}1\\ 3\end{pmatrix}
=\sigma_1\bu_1,
\qquad
A\bv_2=\frac1{\sqrt2}\begin{pmatrix}-3\\ 1\end{pmatrix}
=\sqrt{5}\cdot\frac1{\sqrt{10}}\begin{pmatrix}-3\\ 1\end{pmatrix}
=\sigma_2\bu_2 .
\]
La división por $\sqrt{10}$ deja a $\bu_1$ y $\bu_2$ ortonormales. En
resumen,
\[
U=\frac1{\sqrt{10}}\begin{bmatrix}1&-3\\ 3&1\end{bmatrix},
\qquad
\Sigma=\begin{bmatrix}\sqrt{45}&\\ &\sqrt5\end{bmatrix},
\qquad
V=\frac1{\sqrt2}\begin{bmatrix}1&-1\\ 1&1\end{bmatrix} .
\]
Y la descomposición en sumandos de rango uno, con
$\sqrt2\sqrt{10}=\sqrt{20}$:
\[
\sigma_1\bu_1\bv_1^T+\sigma_2\bu_2\bv_2^T
=\frac{\sqrt{45}}{\sqrt{20}}\begin{bmatrix}1&1\\ 3&3\end{bmatrix}
+\frac{\sqrt{5}}{\sqrt{20}}\begin{bmatrix}3&-3\\ -1&1\end{bmatrix}
=\frac32\begin{bmatrix}1&1\\ 3&3\end{bmatrix}
+\frac12\begin{bmatrix}3&-3\\ -1&1\end{bmatrix}
=\begin{bmatrix}3&0\\ 4&5\end{bmatrix}=A .
\]
\end{example}

\subsection{Algunos casos particulares}

\begin{example}[Una matriz sdp]
Si $S=Q\Lambda Q^T$ es simétrica y definida positiva, su SVD es esa misma
descomposición: $U=V=Q$ y $\Sigma=\Lambda$. Para una matriz sdp, valores
singulares y autovalores coinciden.
\end{example}

\begin{example}[Una simétrica con autovalores negativos]
Si $S=Q\Lambda Q^T$ es simétrica y $S\bx=-\alpha\bx$ con $\alpha>0$, el valor
singular correspondiente es $\sigma=+\alpha$, y uno de los dos vectores
singulares (el $\bu$ o el $\bv$) hay que tomarlo con el signo cambiado: con
$\bv=\bx$ y $\bu=-\bx$, la igualdad $S\bx=-\alpha\bx$ es exactamente
$S\bv=\sigma\bu$. En general, para $S$ simétrica,
$\sigma_k=|\lambda_k|$ reordenados de mayor a menor.
\end{example}

\begin{example}[Una matriz ortogonal]
Si $A=Q$ es ortogonal, $A^TA=I$ y todos los valores singulares valen $1$,
o sea $\Sigma=I$. Los vectores singulares no son únicos: $Q=QII^T$ sirve,
pero también $Q=(QQ_1)IQ_1^T$ para cualquier $Q_1$ ortogonal. Esto es
coherente con $\|Q\|_2=1$: una matriz ortogonal no estira nada.
\end{example}

% ===================================================================
\section{Demostración de la SVD}
\label{sec:prueba-SVD}

Queda por verificar lo que dimos por sentado: que los $\bu_k$ de la
Definición~\ref{def:u} son ortonormales, y que se los puede completar hasta
una base ortonormal de $\R^m$.

\begin{proof}[Demostración del Teorema~\ref{teo:SVD}]
\emph{Los $\bu_k$ son ortonormales.} Para $j,k\le r$,
\[
\bu_j^T\bu_k
=\Bigl(\frac{A\bv_j}{\sigma_j}\Bigr)^T\Bigl(\frac{A\bv_k}{\sigma_k}\Bigr)
=\frac{\bv_j^T(A^TA\,\bv_k)}{\sigma_j\sigma_k}
=\frac{\lambda_k}{\sigma_j\sigma_k}\,\bv_j^T\bv_k
=\frac{\sigma_k}{\sigma_j}\,\bv_j^T\bv_k
=\begin{cases}1&\text{si }j=k,\\ 0&\text{si }j\neq k,\end{cases}
\]
usando que los $\bv$ son ortonormales y que $\lambda_k=\sigma_k^2$.

\emph{Los $\bu_k$ son autovectores de $AA^T$.} No hace falta para la prueba,
pero lo confirma:
\[
AA^T\bu_k=AA^T\Bigl(\frac{A\bv_k}{\sigma_k}\Bigr)
=A\Bigl(\frac{A^TA\,\bv_k}{\sigma_k}\Bigr)
=A\,\frac{\sigma_k^2\bv_k}{\sigma_k}=\sigma_k^2\bu_k .
\]
El paso decisivo fue $(AA^T)A=A(A^TA)$: mover los paréntesis, que es la ley
asociativa y es la clave de una cantidad enorme de pruebas del álgebra lineal.

\emph{Completar las bases.} Los $\bu_1,\dots,\bu_r$ están en $C(A)$ y son $r$
vectores ortonormales, así que son una base ortonormal de $C(A)$, que tiene
dimensión $r$. Elegimos $\bu_{r+1},\dots,\bu_m$ como cualquier base
ortonormal de $N(A^T)$. Por el Corolario~\ref{coro:cuatro},
$C(A)\perp N(A^T)$ y las dimensiones suman $m$, así que la lista completa es
una base ortonormal de $\R^m$ y $U$ es ortogonal. Lo mismo ya estaba hecho del
lado de $V$: los $\bv_{r+1},\dots,\bv_n$ generan $N(A)$, que es ortogonal a
$C(A^T)=\gen(\bv_1,\dots,\bv_r)$.

Con eso, $AV=U\Sigma$ columna por columna es exactamente
\eqref{eq:Av-sigma-u}, y multiplicando por $V^T$ queda $A=U\Sigma V^T$.
\end{proof}

\begin{remark}[Por qué los $\bu$ se definen y no se eligen]
Podría parecer más simétrico tomar los $\bv$ como autovectores ortonormales
de $A^TA$ y los $\bu$ como autovectores ortonormales de $AA^T$, por separado.
No funciona, por dos razones. Primero, los signos: si $\bu$ es autovector,
$-\bu$ también, y la SVD exige el signo que hace $A\bv_k=\sigma_k\bu_k$ con
$\sigma_k>0$. Segundo, y más serio: cuando un $\sigma$ está repetido hay un
plano entero (o más) de autovectores donde elegir, y una elección
independiente de los dos lados no tiene por qué ser compatible. Definir
$\bu_k=A\bv_k/\sigma_k$ resuelve las dos cosas de un saque. Es exactamente el
punto que hace fallar a la ``prueba rápida'' del Ejercicio~\ref{ej:prueba-rapida}.
\end{remark}

\begin{entrega}{Una construcción de la SVD}\label{ent:construccion-SVD}
La demostración anterior parte de los autovectores de $A^TA$ y construye los
$\bu$. Esta entrega hace el camino inverso --- parte de $AA^T$ y construye los
$\bv$ --- y arma la factorización paso a paso.

Sea $A\in\R^{m\times n}$.
\begin{enumerate}[label=\textbf{(\alph*)},itemsep=3pt,topsep=4pt]
\item Probar que $A^TA$ y $AA^T$ son simétricas, semidefinidas positivas, y
que tienen los mismos autovalores no nulos. \emph{Sugerencia:} para lo último,
si $A^TA\,\bv=\lambda\bv$ con $\lambda\neq0$, ver que $A\bv$ es autovector de
$AA^T$ con el mismo autovalor, y verificar que $A\bv\neq\bi 0$.
\end{enumerate}

\medskip
\noindent Sean $\sigma_1\ge\cdots\ge\sigma_r>0$ tales que
$\sigma_1^2,\dots,\sigma_r^2$ son los autovalores no nulos de $AA^T$, con
$r=\rg(A)$. Sea $U=[\,\bu_1\ \cdots\ \bu_m\,]$ ortogonal con columnas
autovectores de $AA^T$, y $U_r=[\,\bu_1\ \cdots\ \bu_r\,]$. Definimos
$\Sigma_r=\diag(\sigma_1,\dots,\sigma_r)$ y
\[
V_r=A^TU_r\Sigma_r^{-1}\in\R^{n\times r} .
\]
Elegimos además $V_{-r}=[\,\bv_{r+1}\ \cdots\ \bv_n\,]$ una base ortonormal de
$N(A)$, y ponemos $V=[\,V_r\ \ V_{-r}\,]$, $U=[\,U_r\ \ U_{-r}\,]$ y
\[
\Sigma=\begin{bmatrix}\Sigma_r&0\\ 0&0\end{bmatrix}\in\R^{m\times n} .
\]

\begin{enumerate}[label=\textbf{(\alph*)},itemsep=3pt,topsep=4pt,start=2]
\item Probar que $AV_r=U_r\Sigma_r$.
\item Probar que $AV_{-r}=0$.
\item Concluir que $AV=[\,U_r\Sigma_r\ \ 0\,]=U\Sigma$.
\item Probar que $V^TV=I_n$. \emph{Sugerencia:} verificar por separado
$V_r^TV_r=I_r$, $V_r^TV_{-r}=0$ y $V_{-r}^TV_{-r}=I_{n-r}$.
\item Deducir, por ser $V$ cuadrada, que $VV^T=I_n$ y $V^T=V^{-1}$.
\item Concluir que $A=U\Sigma V^T$.
\item Comparar con la construcción de la Sección~\ref{sec:prueba-SVD}: ¿en
qué punto exacto de esta construcción se usa que los $\bu$ ya estaban
elegidos, y por qué eso determina los signos de los $\bv$?
\end{enumerate}

\medskip
\noindent{\footnotesize\emph{Se entrega por escrito. El ítem (a) es el que se
usa en todo el capítulo; el (h) es la contracara de la observación anterior
sobre por qué una de las dos familias se define a partir de la otra.}}
\end{entrega}

% ===================================================================
\section{Lo que se lee en la SVD}
\label{sec:cuatro-SVD}

\subsection{Los cuatro subespacios fundamentales}

El Capítulo~\ref{cap:algebra-lineal} presentó los cuatro subespacios
fundamentales y el Capítulo~\ref{cap:ortogonalidad} probó cómo se aparean
ortogonalmente. La SVD hace algo más: da bases \emph{ortonormales} de los
cuatro al mismo tiempo.

\begin{proposition}\label{prop:cuatro-SVD}
Sea $A=U\Sigma V^T$ con $\rg(A)=r$. Entonces
\begin{center}
\renewcommand{\arraystretch}{1.25}
\begin{tabular}{@{}lll@{}}
\hline
subespacio & base ortonormal & dimensión\\
\hline
espacio fila $C(A^T)$ & $\bv_1,\dots,\bv_r$ & $r$\\
espacio nulo $N(A)$ & $\bv_{r+1},\dots,\bv_n$ & $n-r$\\
espacio columna $C(A)$ & $\bu_1,\dots,\bu_r$ & $r$\\
espacio nulo a izquierda $N(A^T)$ & $\bu_{r+1},\dots,\bu_m$ & $m-r$\\
\hline
\end{tabular}
\end{center}
Además, las proyecciones ortogonales sobre esos subespacios son
\[
P_{C(A)}=U_rU_r^T,
\qquad
P_{C(A^T)}=V_rV_r^T,
\]
y $P_{N(A^T)}=I-U_rU_r^T$, $P_{N(A)}=I-V_rV_r^T$.
\end{proposition}

\begin{proof}
De \eqref{eq:Av-sigma-u}, $A\bv_k=\bi 0$ para $k>r$, así que
$\bv_{r+1},\dots,\bv_n$ son $n-r$ vectores ortonormales de $N(A)$, que tiene
dimensión $n-r$: son una base. Los $\bv_1,\dots,\bv_r$ son ortonormales y
ortogonales a $N(A)$, luego están en $N(A)^\perp=C(A^T)$
(Corolario~\ref{coro:cuatro}), y son $r$: base. Del lado de $U$, los
$\bu_k=A\bv_k/\sigma_k$ con $k\le r$ están en $C(A)$ y son $r$ ortonormales:
base de $C(A)$; los restantes completan y están en $C(A)^\perp=N(A^T)$.

Para las proyecciones: $U_r$ tiene columnas ortonormales que generan $C(A)$,
y ya vimos en la Sección~\ref{sec:proyecciones} que la proyección ortogonal
sobre el espacio generado por columnas ortonormales es $QQ^T$.
\end{proof}

\begin{remark}
Vale la pena leer la tabla junto con \eqref{eq:Av-sigma-u}: la SVD dice que
$A$ manda una base ortonormal del espacio fila a una base ortogonal del
espacio columna, estirando cada dirección por su $\sigma_k$, y aplasta el
espacio nulo a cero. Ésa es toda la información que hay en una matriz.
\end{remark}

\subsection{Los vectores singulares de $A^T$}

Trasponiendo $A=U\Sigma V^T$ se obtiene $A^T=V\Sigma^TU^T$, que recorre el
camino inverso: de los $\bu$ a los $\bv$.

\begin{proposition}\label{prop:AT}
$A^T\bu_k=\sigma_k\bv_k$ para $k=1,\dots,r$, y $A^T\bu_k=\bi 0$ para
$k=r+1,\dots,m$. En particular $A$ y $A^T$ tienen los mismos valores
singulares no nulos, y $\|A\|_2=\|A^T\|_2$.
\end{proposition}

\begin{proof}
Multiplicando $A\bv_k=\sigma_k\bu_k$ por $A^T$ y usando
$A^TA\bv_k=\sigma_k^2\bv_k$ queda $\sigma_k^2\bv_k=\sigma_kA^T\bu_k$;
dividiendo por $\sigma_k\neq0$ sale la igualdad. Para $k>r$, $\bu_k\in N(A^T)$
por la Proposición~\ref{prop:cuatro-SVD}.
\end{proof}

\subsection{Una matriz simétrica que también da la SVD}

Construimos la SVD a partir de $A^TA$ y $AA^T$. Hay otra manera, que usa una
sola matriz simétrica de tamaño $m+n$ y tiene $r$ pares de autovalores
opuestos:
\[
S=\begin{bmatrix}0&A\\ A^T&0\end{bmatrix}
\qquad\text{tiene autovectores}\qquad
\begin{pmatrix}\bu_k\\ \bv_k\end{pmatrix}
\quad\text{y}\quad
\begin{pmatrix}-\bu_k\\ \bv_k\end{pmatrix} .
\]
En efecto, usando $A\bv_k=\sigma_k\bu_k$ y $A^T\bu_k=\sigma_k\bv_k$,
\begin{equation}\label{eq:bloques-SVD}
\begin{bmatrix}0&A\\ A^T&0\end{bmatrix}
\begin{pmatrix}\pm\bu_k\\ \bv_k\end{pmatrix}
=\begin{pmatrix}A\bv_k\\ \pm A^T\bu_k\end{pmatrix}
=\sigma_k\begin{pmatrix}\bu_k\\ \bv_k\end{pmatrix}
\quad\text{y}\quad
-\sigma_k\begin{pmatrix}-\bu_k\\ \bv_k\end{pmatrix} .
\end{equation}
Eso da $2r$ autovalores, $\pm\sigma_1,\dots,\pm\sigma_r$. Los autovectores son
ortogonales: $(-\bu_k)^T\bu_k+\bv_k^T\bv_k=-1+1=0$. Los otros
$(m-r)+(n-r)$ autovectores, correspondientes a $\lambda=0$, se arman con los
$\bu$ y $\bv$ que quedaron en $N(A^T)$ y $N(A)$
(Ejercicio~\ref{ej:bloques}).

\subsection{$A^TA$ y $AA^T$: los mismos valores singulares}

La construcción partió de $A^TA$, pero podríamos haber empezado por $AA^T$ y
nada habría cambiado. La razón es un hecho general sobre autovalores, que no
tiene nada que ver con la SVD y que por eso probamos en su lugar: la
Proposición~\ref{prop:AB-BA} del Capítulo~\ref{cap:autovalores} dice que $AB$
y $BA$ tienen los mismos autovalores no nulos, aun siendo rectangulares.

Con $B=A^T$ eso dice que $A^TA$ y $AA^T$ tienen los mismos autovalores no
nulos, o sea que $A$ tiene \emph{los mismos valores singulares} se los calcule
como se los calcule. Lo que cambia es sólo el tamaño: si $m>n$, la matriz
$AA^T$ es de $m\times m$ y tiene $m-n$ autovalores \emph{cero} de más respecto
de $A^TA$, que es de $n\times n$. Eso explica la asimetría de la
Figura~\ref{fig:esquema-SVD}: sobran columnas en $U$, no en $V$. En la
práctica conviene diagonalizar la más chica de las dos, y es exactamente lo
que hace el truco de la matriz de Gram del capítulo de PCA.

% ===================================================================
\section{La geometría: de la esfera al elipsoide}
\label{sec:geometria-SVD}

La SVD factoriza $A$ como (ortogonal) $\times$ (diagonal) $\times$
(ortogonal), y esos tres factores tienen una lectura geométrica inmediata:
$V^T$ y $U$ rotan (posiblemente con una reflexión) y $\Sigma$ estira a lo
largo de los ejes. En dos dimensiones se puede dibujar todo el proceso
(Figura~\ref{fig:svd-geometria}): un vector $\bx$ de la circunferencia
unitaria se rota a $V^T\bx$, se estira a $\Sigma V^T\bx$ y se vuelve a rotar
hasta $A\bx=U\Sigma V^T\bx$, que recorre una elipse.

\begin{figure}[ht]
\centering
\begin{tikzpicture}[scale=0.98]
\draw (1cm,3cm) circle (0.7cm);%first circle
\draw (0cm,3cm)--(2cm,3cm);
\draw (1cm,4cm)--(1cm,2cm);
\draw[very thick,->] (1cm,3cm)--(1.3cm,3.64cm);
\node at (1.4cm,3.8cm) {${\bv_2}$};
\draw[very thick,->] (1cm,3cm)--(1.64cm,2.7cm);
\node at (1.9cm,2.7cm) {${\bv_1}$};
\draw [<-](2.8cm,4.0cm) arc [radius=1cm, start angle=50, end angle=125];
\node at (2.2cm,4.5cm) {\color{bleu}${V^T}$};
\draw [<-](1.5cm,2.0cm) arc [radius=1cm, start angle=230, end angle=310];
\node at (2.15cm,1.55cm) {${V}$};
\node at(1.5cm,4.9cm){$\bx$};
\node at(10.6cm,4.9cm){$A\bx$};
\draw (3.5cm,3cm) circle (0.7cm);%second circle
\draw (2.5cm,3cm)--(4.5cm,3cm);
\draw (3.5cm,4cm)--(3.5cm,2cm);
\draw [very thick,->] (3.5cm,3cm)--(3.5cm,3.7cm);
\draw [very thick,->] (3.5cm,3cm)--(4.2cm,3cm);
\draw [<-](5.5,4.0cm) arc [radius=1cm, start angle=50, end angle=125];
\node at (4.9,4.5cm) {\color{bleu}${\Sigma}$};
\draw[rotate=0] (7cm, 3cm) ellipse (1.6cm and 0.4cm);
\draw (5.1cm,3cm)--(8.9cm,3cm);
\draw [very thick,->] (7cm,3cm)--(8.6cm,3cm);
\node at (8.8cm,3.2cm) {${\sigma_1}$};
\draw (7cm, 4.2cm)--(7cm, 1.8cm);
\draw [very thick,->] (7cm, 3cm)--(7cm, 3.4cm);
\node at (7.3cm, 3.6cm) {${\sigma_2}$};
\draw[rotate=135] (-5.8cm, -10.04cm) ellipse (1.6cm and 0.4cm);
\draw (9.6cm,3cm)--(12.9cm,3cm);
\draw (11.2cm, 4.2cm)--(11.2cm, 1.8cm);
\draw [very thick,->] (11.2cm, 3cm)--(12.33,1.87);
\draw [dash pattern=on 2pt off 2pt] (11.2cm,3cm)--(9.9cm,4.3cm);
\draw [very thick,->] (11.2cm,3cm)--(11.48cm,3.28cm);
\draw [dash pattern=on 2pt off 2pt] (11.2cm,3cm)--(10.7cm,2.5cm);
\draw [<-](9.5,4.0cm) arc [radius=1cm, start angle=50, end angle=125];
\node at (8.9,4.5cm) {\color{bleu}${U}$};
\node at (12.43,1.6) {${\sigma_1\bu_1}$};
\node at (11.98cm,3.3cm) {${\sigma_2\bu_2}$};
\draw [<-,dash pattern=on 2pt off 2pt](10.2,4.9cm) arc [radius=10cm, start angle=65, end angle=115];
\node at (5.78cm,6.1cm) {\color{bleu}${A}$};
\end{tikzpicture}
\caption{Rotación, estiramiento y rotación. $V^T$ y $U$ son rotaciones
(posiblemente con reflexión) y $\Sigma$ estira los ejes por $\sigma_1$ y
$\sigma_2$. La circunferencia unitaria va a parar a una elipse de semiejes
$\sigma_1$ y $\sigma_2$ en las direcciones $\bu_1$ y $\bu_2$.}
\label{fig:svd-geometria}
\end{figure}

El enunciado preciso es el siguiente, y conecta con la
Sección~\ref{sec:elipses}, donde estudiamos las elipses asociadas a una
matriz sdp.

\begin{theorem}[La imagen de la esfera]\label{teo:elipsoide}
Sea $A\in\R^{n\times n}$ invertible, con valores singulares
$\sigma_1\ge\cdots\ge\sigma_n>0$. Entonces la imagen de la esfera unitaria,
\[
\{A\bx : \|\bx\|=1\} ,
\]
es el elipsoide de semiejes $\sigma_1,\dots,\sigma_n$ en las direcciones
$\bu_1,\dots,\bu_n$. En particular, $\sigma_1$ y $\sigma_n$ son el semieje
mayor y el semieje menor: $A$ estira a lo sumo $\sigma_1$ y a lo menos
$\sigma_n$.
\end{theorem}

\begin{proof}
Sea $\by=A\bx$ con $\|\bx\|=1$. Escribiendo $\bx=\sum_k c_k\bv_k$, es
$\by=\sum_k c_k\sigma_k\bu_k$, y las coordenadas de $\by$ en la base
ortonormal $\{\bu_k\}$ son $y'_k=\bu_k^T\by=\sigma_kc_k$. La condición
$\|\bx\|^2=\sum c_k^2=1$ se traduce entonces en
\[
\sum_{k=1}^n\frac{(y'_k)^2}{\sigma_k^2}=1 ,
\]
que es la ecuación de un elipsoide de semiejes $\sigma_k$ en las direcciones
$\bu_k$. Recíprocamente, todo punto que cumple esa ecuación es la imagen del
$\bx$ correspondiente.
\end{proof}

\begin{remark}[Es la elipse del Capítulo~\ref{cap:sdp}]
El elipsoide del teorema es exactamente el conjunto
$\{\by : \by^T(AA^T)^{-1}\by=1\}$, o sea la elipse asociada a la matriz sdp
$(AA^T)^{-1}$ en el sentido de la Sección~\ref{sec:elipses}: sus ejes son los
autovectores $\bu_k$ de $AA^T$ y sus semiejes son
$1/\sqrt{\lambda_k((AA^T)^{-1})}=\sqrt{\lambda_k(AA^T)}=\sigma_k$. La SVD es,
leída así, el teorema que dice que \emph{toda} matriz manda esferas a
elipsoides, y que los ejes de llegada son los $\bu_k$ y los de partida los
$\bv_k$.
\end{remark}

\begin{remark}[Contar parámetros]
La factorización ordena la información de una matriz de $2\times2$: los cuatro
números $a,b,c,d$ se corresponden con dos ángulos $\theta$, $\phi$ y dos
valores singulares,
\[
\begin{bmatrix}a&b\\ c&d\end{bmatrix}
=\begin{bmatrix}\cos\theta&-\sen\theta\\ \sen\theta&\cos\theta\end{bmatrix}
\begin{bmatrix}\sigma_1&\\ &\sigma_2\end{bmatrix}
\begin{bmatrix}\cos\phi&\sen\phi\\ -\sen\phi&\cos\phi\end{bmatrix} .
\]
Si la matriz es simétrica hay sólo tres parámetros independientes: ¿cómo se
reducen entonces los cuatro? El Ejercicio~\ref{ej:parametros} sigue esta
cuenta en dimensiones más altas.
\end{remark}

% ===================================================================
\section{La descomposición polar}
\label{sec:polar}

Todo número complejo se escribe en forma polar como $z=re^{i\theta}$, con
$r\ge0$ y $|e^{i\theta}|=1$. Pensando a los números como matrices de
$1\times1$, el factor $e^{i\theta}$ hace de matriz ortogonal y $r$ de matriz
semidefinida positiva. La SVD dice que esa descomposición existe en cualquier
dimensión.

\begin{theorem}[Descomposición polar]\label{teo:polar}
Toda $A\in\R^{n\times n}$ se factoriza como
\[
\boxed{\ A=QS\ }
\]
con $Q$ ortogonal y $S$ simétrica semidefinida positiva. Si $A$ es
invertible, $S$ es definida positiva y $S=(A^TA)^{1/2}$ es la raíz cuadrada
de la Sección~\ref{sec:raiz}.
\end{theorem}

\begin{proof}
Sea $A=U\Sigma V^T$. Intercalando $V^TV=I$,
\[
A=U\Sigma V^T=(UV^T)(V\Sigma V^T)=QS,
\qquad Q:=UV^T,\quad S:=V\Sigma V^T .
\]
$Q$ es ortogonal por ser producto de ortogonales. $S$ es simétrica y sus
autovalores son los $\sigma_k\ge0$, luego es semidefinida positiva, y
definida positiva si todos los $\sigma_k>0$, es decir si $A$ es invertible.
Finalmente $S^2=V\Sigma^2V^T=V\Sigma^TU^TU\Sigma V^T=A^TA$, y como $S$ es
sdp, es \emph{la} raíz cuadrada de $A^TA$ por la unicidad probada en la
Entrega~\ref{ent:raiz}.
\end{proof}

También vale la versión con los factores al revés, $A=S_1Q$ con
$S_1=U\Sigma U^T=(AA^T)^{1/2}$ y el mismo $Q=UV^T$.

\begin{example}
Para $A=\begin{bsmallmatrix}3&0\\ 4&5\end{bsmallmatrix}$, con la SVD del
Ejemplo~\ref{ej:svd-completo},
\[
Q=UV^T=\frac1{\sqrt5}\begin{bmatrix}2&-1\\ 1&2\end{bmatrix},
\qquad
S=V\Sigma V^T=\sqrt5\begin{bmatrix}2&1\\ 1&2\end{bmatrix} .
\]
\end{example}

\begin{remark}[Rotación y deformación]
En mecánica del continuo la descomposición polar separa la \emph{rotación} de
un cuerpo (el factor $Q$) de su \emph{deformación} (el factor $S$): los
autovalores de $S$ son los factores de estiramiento y sus autovectores, las
direcciones principales de deformación. Es la misma lectura que la del
Teorema~\ref{teo:elipsoide}, reagrupada.
\end{remark}

% ===================================================================
\section{Las tres normas de la SVD}
\label{sec:tres-normas}

La norma de Frobenius apareció en la Sección~\ref{sec:frobenius} y la
espectral en la Sección~\ref{sec:sigma1}. Hasta ahora no tenían nada que ver
una con otra: la primera lee la matriz como un vector largo y la segunda mide
cuánto estira. La SVD las pone a las dos ---y a una tercera--- en el mismo
idioma: \emph{todas se leen de los valores singulares}. Esta sección es corta
y es la que van a citar los dos capítulos siguientes.

El ingrediente es que multiplicar por una matriz de columnas ortonormales no
cambia la norma de Frobenius: si $Q^TQ=I$, entonces
\[
\|QA\|_F^2=\tr\bigl((QA)^T(QA)\bigr)=\tr(A^TQ^TQA)=\tr(A^TA)=\|A\|_F^2 ,
\]
y lo mismo del otro lado, porque $\|AP\|_F^2=\tr(APP^TA^T)$ y
$\|A\|_F^2=\tr(AA^T)$ también.

\begin{proposition}[Frobenius cuenta todos los valores singulares]
\label{prop:frobenius-sigmas}
Sea $A\in\R^{m\times n}$ con valores singulares
$\sigma_1\ge\cdots\ge\sigma_r>0$. Entonces
\[
\boxed{\ \|A\|_F^2=\sigma_1^2+\sigma_2^2+\cdots+\sigma_r^2\ }
\]
\end{proposition}

\begin{proof}
Con la SVD reducida $A=U_r\Sigma_rV_r^T$ de~\eqref{eq:SVD-reducida}, las
columnas de $U_r$ y las de $V_r$ son ortonormales, así que
\[
\|A\|_F=\|U_r\Sigma_rV_r^T\|_F=\|\Sigma_rV_r^T\|_F=\|\Sigma_r\|_F ,
\]
y $\Sigma_r=\diag(\sigma_1,\dots,\sigma_r)$ tiene norma de Frobenius
$\bigl(\sum_k\sigma_k^2\bigr)^{1/2}$.
\end{proof}

Comparando con $\|A\|_2=\sigma_1$ (Teorema~\ref{teo:sigma1}) sale de una
\[
\|A\|_2\ \le\ \|A\|_F\ \le\ \sqrt{r}\,\|A\|_2 ,
\]
con la primera igualdad exactamente cuando $A$ tiene rango uno y la segunda
cuando todos los valores singulares no nulos coinciden.

Hay una tercera manera obvia de resumir los $\sigma_k$ en un número, y
también es una norma.

\begin{definition}[Norma nuclear]\label{def:norma-nuclear}
La \textbf{norma nuclear} (o \emph{norma de traza}) de $A$ es la suma de sus
valores singulares,
\[
\|A\|_*=\sigma_1+\sigma_2+\cdots+\sigma_r .
\]
\end{definition}

Que sea una norma no es evidente desde esta definición ---la desigualdad
triangular dice que los valores singulares de $A+B$ suman menos que los de
$A$ más los de $B$, y eso hay que probarlo--- y lo dejamos para el capítulo
de normas de matrices. Las tres juntas:
\[
\|A\|_2=\sigma_1,\qquad
\|A\|_F=\bigl(\sigma_1^2+\cdots+\sigma_r^2\bigr)^{1/2},\qquad
\|A\|_*=\sigma_1+\cdots+\sigma_r ,
\]
y son genuinamente distintas, ya que para la identidad de $n\times n$ dan
$1$, $\sqrt n$ y $n$.

Lo que tienen en común es más importante que lo que las diferencia.

\begin{definition}[Norma ortogonalmente invariante]\label{def:ortog-invariante}
Una norma $\|\cdot\|$ en $\R^{m\times n}$ es \textbf{ortogonalmente
invariante} si
\[
\|QAP\|=\|A\| \qquad\text{para toda } A
\]
y todas las matrices ortogonales $Q\in\R^{m\times m}$, $P\in\R^{n\times n}$.
\end{definition}

\begin{corollary}\label{coro:tres-invariantes}
Las normas espectral, de Frobenius y nuclear son ortogonalmente invariantes.
\end{corollary}

\begin{proof}
Si $A=U\Sigma V^T$, entonces $QAP=(QU)\Sigma(P^TV)^T$ es una SVD de $QAP$,
porque $QU$ y $P^TV$ son ortogonales. Luego $QAP$ tiene \emph{los mismos
valores singulares} que $A$, y las tres normas se calculan con ellos.
\end{proof}

Esa es la propiedad que hace que la SVD sea la herramienta correcta para
medir distancias entre matrices: si una norma no ve más que los valores
singulares, la SVD ya la calculó. El capítulo de PCA se apoya exactamente en
esto para probar que la SVD truncada es la mejor aproximación de rango bajo.

\begin{remark}[Cuál es cuál]
Conviene quedarse con una imagen de cada una. La espectral mide el
\emph{peor caso}: cuánto puede estirar $A$ a un vector. La de Frobenius mide
la \emph{energía total}, repartida entre todas las direcciones. La nuclear
es la que mejor detecta el \emph{rango}: es la suma de los $\sigma_k$, así
que penaliza tener muchos valores singulares chicos y no solo uno grande ---por
eso aparece cuando se quiere forzar que una matriz tenga rango bajo.
\end{remark}

% ===================================================================
\section{Lo que viene}
\label{sec:lo-que-viene}

Cerramos enunciando, sin demostrar, los dos resultados que esta factorización
hace posibles y que se tratan en los capítulos siguientes.

\medskip
\noindent\textbf{Aproximación de rango bajo.} Truncar la
suma~\eqref{eq:sumandos} en el término $d$-ésimo,
\[
A_d=\sigma_1\bu_1\bv_1^T+\cdots+\sigma_d\bu_d\bv_d^T ,
\]
da una matriz de rango $d$. El \textbf{teorema de Eckart-Young} afirma que es
la mejor posible:
\[
\text{si } \rg(B)\le d, \text{ entonces } \|A-A_d\|\le\|A-B\| ,
\]
tanto en la norma de Frobenius como en la espectral. Ésa es la razón por la
que la SVD truncada se usa para comprimir imágenes y para reducir la
dimensión de un conjunto de datos. El enunciado preciso, la demostración y la
lectura estadística son el contenido del capítulo de PCA.

\medskip
\noindent\textbf{La pseudoinversa.} Si $A$ no es invertible --- o ni siquiera
cuadrada --- no hay $A^{-1}$, pero la SVD permite hacer lo más parecido
posible: invertir lo que se puede invertir y dejar en cero el resto. Con
$\Sigma^+\in\R^{n\times m}$ la matriz que tiene $1/\sigma_1,\dots,1/\sigma_r$
en la diagonal y ceros en el resto, se define
\[
A^+=V\,\Sigma^+\,U^T .
\]
Se verá que $AA^+$ y $A^+A$ son las proyecciones ortogonales sobre $C(A)$ y
$C(A^T)$ --- comparar con la Proposición~\ref{prop:cuatro-SVD} ---, que
$A^+=(A^TA)^{-1}A^T$ cuando las columnas de $A$ son independientes, y que
$A^+\bb$ es la solución de cuadrados mínimos de norma mínima. Es el contenido
del capítulo final.

% ===================================================================
\section*{Ejercicios de desarrollo}

\begin{entrega}{La norma espectral}\label{ent:norma-espectral}
Todo este ejercicio se apoya en la Definición~\ref{def:norma-espectral} y en
el Teorema~\ref{teo:sigma1}; casi nada necesita la SVD.

\begin{enumerate}[label=\textbf{(\alph*)},itemsep=3pt,topsep=4pt]
\item Verificar directamente sobre la definición que
\[
\|A\|_2=\max_{\|\bx\|=1}\|A\bx\|=\max_{0<\|\bx\|\le1}\|A\bx\| ,
\]
es decir, que las tres expresiones dan el mismo número. (Que el máximo se
alcance, y no haya que escribir supremo, ya lo sabemos: el
Teorema~\ref{teo:sigma1} exhibe el vector donde se alcanza.)

\item Calcular $\|A\|_2$ para
$A=\begin{bsmallmatrix}1&0\\ 0&3\end{bsmallmatrix}$,
$A=\begin{bsmallmatrix}0&2\\ 0&0\end{bsmallmatrix}$ y
$A=\begin{bsmallmatrix}1&1\\ 0&1\end{bsmallmatrix}$. En el último caso,
comparar $\|A\|_2$ con $\max_i|\lambda_i(A)|$ y observar que la desigualdad
del Corolario~\ref{coro:lambda-sigma} es estricta.

\item Probar que $\|A\|_2=\|A^T\|_2$. \emph{Sugerencia:} $A^TA$ y $AA^T$
tienen los mismos autovalores no nulos
(Proposición~\ref{prop:AB-BA}). Mostrar en cambio con un ejemplo que
\emph{no} es cierto que $\|A\bx\|=\|A^T\bx\|$ para todo $\bx$, y probar que
sí lo es cuando $A$ conmuta con su traspuesta.

\item Probar que $\|QA\|_2=\|A\|_2=\|AQ'\|_2$ si $Q$ y $Q'$ son ortogonales
de tamaños compatibles. \emph{Sugerencia:} las ortogonales no cambian la
norma de los vectores.

\item \emph{(Una submatriz no puede estirar más que la matriz entera.)} Sea
$B$ la submatriz de $A$ formada por algunas de sus filas y algunas de sus
columnas. Probar que $\sigma_1(B)\le\sigma_1(A)$, es decir
$\|B\|_2\le\|A\|_2$, de las dos maneras siguientes.
\begin{enumerate}[label=(\roman*)]
\item \emph{A mano.} Si $C$ se obtiene de $A$ quedándose con algunas
columnas y $\by$ es un vector del tamaño correspondiente, sea $\bx\in\R^n$ el
vector que resulta de rellenar con ceros las posiciones eliminadas: verificar
que $A\bx=C\by$ y $\|\bx\|=\|\by\|$, y concluir $\|C\|_2\le\|A\|_2$. Después,
si $D$ sale de $C$ quedándose con algunas filas, observar que $D\by$ es
$C\by$ con coordenadas borradas, así que $\|D\by\|\le\|C\by\|$. Encadenar las
dos operaciones.
\item \emph{En forma matricial.} Escribir $B=R\,A\,C^T$ con $R$ y $C$
matrices de ceros y unos que seleccionan filas y columnas, verificar que $R$ y
$C$ tienen filas ortonormales y deducir la desigualdad del ítem (d).
\end{enumerate}

\item Sea $A$ con valores singulares $\sigma_1,\dots,\sigma_r$. Determinar
los valores singulares de $A^T$, de $A^TA$, de $AA^TA$, y de $A^{-1}$ cuando
$A$ es cuadrada e invertible. ¿Cuánto vale $\|A^{-1}\|_2$?

\item \emph{(La norma espectral no viene de un producto interno.)} Tomar
$A=\begin{bsmallmatrix}1&0\\ 0&0\end{bsmallmatrix}$ y
$B=\begin{bsmallmatrix}0&0\\ 0&1\end{bsmallmatrix}$. Calcular
$\|A+B\|_2$, $\|A-B\|_2$, $\|A\|_2$, $\|B\|_2$ y verificar que falla la ley
del paralelogramo. Concluir, por el Teorema~\ref{teo-normaPI-carac}, que
$\|\cdot\|_2$ no proviene de ningún producto interno. Repetir la cuenta con
$\|\cdot\|_F$ y las mismas dos matrices, y explicar por qué esta vez la ley
\emph{tiene} que cumplirse.
\end{enumerate}

\medskip
\noindent{\footnotesize\emph{Se entrega por escrito. El ítem (b) muestra que
$\sigma_1$ y $|\lambda|_{\max}$ son cosas distintas incluso para matrices
cuadradas; el (g) es la diferencia de fondo entre las dos normas de matrices
que manejamos.}}
\end{entrega}

\begin{entrega}{La SVD y los cuatro subespacios}\label{ent:cuatro-SVD}
Sea $A\in\R^{m\times n}$ de rango $r$, con SVD $A=U\Sigma V^T$.

\begin{enumerate}[label=\textbf{(\alph*)},itemsep=3pt,topsep=4pt]
\item Calcular la SVD completa de
$A=\begin{bmatrix}0&2&0\\ 0&0&3\\ 0&0&0\end{bmatrix}$
y exhibir las bases ortonormales de los cuatro subespacios que da la
Proposición~\ref{prop:cuatro-SVD}. (Acá $U$ y $V$ son matrices de
permutación.)

\item Verificar sobre esa matriz que $U_rU_r^T$ y $V_rV_r^T$ son las
proyecciones ortogonales sobre $C(A)$ y $C(A^T)$: comprobar $P^2=P=P^T$ y
$\tr(P)=r$.

\item Probar en general que $U_rU_r^TA=A$ y $AV_rV_r^T=A$, e interpretar
las dos igualdades en palabras.

\item Probar que $\sigma_1(A)\ge\sigma_1(A_d)$ y, más precisamente, que los
valores singulares de $A_d=\sum_{k\le d}\sigma_k\bu_k\bv_k^T$ son
$\sigma_1,\dots,\sigma_d$, y los de $A-A_d$ son
$\sigma_{d+1},\dots,\sigma_r$. Deducir $\|A-A_d\|_2=\sigma_{d+1}$.

\item \emph{(La esfera y el elipsoide.)} Para
$A=\begin{bsmallmatrix}3&0\\ 4&5\end{bsmallmatrix}$ del
Ejemplo~\ref{ej:svd-completo}, dibujar la circunferencia unitaria, marcar
$\bv_1$ y $\bv_2$, y dibujar su imagen $\{A\bx:\|\bx\|=1\}$ marcando
$\sigma_1\bu_1$ y $\sigma_2\bu_2$. Verificar sobre el dibujo que
$\bv_1\perp\bv_2$ y que $\bu_1\perp\bu_2$, y observar que $A$ \emph{no}
preserva los ángulos entre otros pares de vectores.

\item Probar que $\{A\bx:\|\bx\|\le1\}$ es el elipsoide \emph{lleno}, y que
su volumen es $|\det A|$ veces el de la bola unitaria. \emph{Sugerencia:}
$|\det A|=\sigma_1\cdots\sigma_n$ (probarlo también).
\end{enumerate}

\medskip
\noindent{\footnotesize\emph{Se entrega por escrito. El ítem (d) es el que se
usa en el capítulo de PCA para medir el error de la SVD truncada; el (f)
explica por qué el determinante mide volúmenes.}}
\end{entrega}

% ===================================================================
\section*{Ejercicios}

\begin{enumerate}

\item Una matriz simétrica $S$ tiene autovectores ortonormales
$\bv_1,\dots,\bv_n$, así que todo $\bx$ se escribe
$\bx=c_1\bv_1+\cdots+c_n\bv_n$. Explicar estas dos fórmulas:
\[
\bx^T\bx=c_1^2+\cdots+c_n^2,
\qquad
\bx^TS\bx=\lambda_1c_1^2+\cdots+\lambda_nc_n^2 .
\]

\item El ejercicio anterior da una forma clara del cociente de Rayleigh:
\[
R(\bx)=\frac{\bx^TS\bx}{\bx^T\bx}
=\frac{\lambda_1c_1^2+\cdots+\lambda_nc_n^2}{c_1^2+\cdots+c_n^2} .
\]
¿Por qué el valor máximo de este cociente es $\lambda_1$? ¿Y por qué se
alcanza cuando $c_1=1$ y $c_2=\cdots=c_n=0$? Ésta es la forma más simple de
entender el Teorema~\ref{teo:sigma1}.

\item Siguiendo el ejercicio anterior: maximizar $R(\bx)$ con la condición
$\bx^T\bv_1=0$. ¿Qué significa esa condición para $c_1$? ¿Por qué el máximo
es ahora $\lambda_2$? ¿Y sujeto a qué dos condiciones sobre $\bx$ el máximo
vale $\lambda_3$?

\item Para maximizar $\bx^TS\bx$ con la restricción $\bx^T\bx=1$, el método
de Lagrange trabaja con $L(\bx)=\bx^TS\bx-\lambda(\bx^T\bx-1)$. Mostrar que
$\nabla L=\bi 0$ es exactamente $S\bx=\lambda\bx$, de modo que otra vez
$\max R(\bx)=\lambda_1$.

\item Encontrar los $\sigma$, los $\bv$ y los $\bu$ en la SVD de
$A=\begin{bsmallmatrix}3&4\\ 0&5\end{bsmallmatrix}$.

\item Encontrar la SVD de la matriz de rango uno
$A=\begin{bsmallmatrix}2&4\\ 1&2\end{bsmallmatrix}$, y factorizar $A^TA$
como $Q\Lambda Q^T$.

\item Sea $A$ con valores singulares $\sigma_1\ge\cdots\ge\sigma_r>0$ y sea
$B=A-\sigma_1\bu_1\bv_1^T$ la matriz que queda al quitarle a $A$ su
sumando de rango uno más grande. ¿Cuáles son los valores singulares de $B$?
¿Cuánto vale $\|B\|_2$? ¿Cuál es el rango de $B$?

\item\label{ej:bloques} Verificar que la traza de
$S=\begin{bsmallmatrix}0&A\\ A^T&0\end{bsmallmatrix}$, sumando sus elementos
diagonales, coincide con la suma de los autovalores de
\eqref{eq:bloques-SVD}. Si $A$ es diagonal de $n\times n$ con elementos
$1,2,\dots,n$, ¿cuáles son los $2n$ autovalores y autovectores de $S$?
Describir además los autovectores correspondientes a $\lambda=0$ en el caso
general.

\item\label{ej:reducida} La ecuación \eqref{eq:SVD-reducida} es
$AV_r=U_r\Sigma_r$. Multiplicando por $V_r^T$ se obtiene
$A=U_r\Sigma_rV_r^T$. Para ese paso \emph{no} se puede usar $V_rV_r^T=I$
(que es falso cuando $r<n$). Mostrar en cambio que la matriz
$U_r\Sigma_rV_r^T$ satisface \eqref{eq:Av-sigma-u}, y concluir que es $A$.

\item\label{ej:prueba-rapida} He aquí una prueba rápida de la SVD. El paso 2
usa las factorizaciones $A^TA=V\Lambda V^T$ y $AA^T=U\Lambda U^T$ (con los
mismos autovalores en $\Lambda$):
\begin{center}
\begin{tabular}{ll@{\quad}l}
\hline
1 & $A(A^TA)$ & $=(AA^T)A$\\
2 & $AV\Lambda V^T$ & $=U\Lambda U^TA$\\
3 & $(U^TAV)\Lambda$ & $=\Lambda(U^TAV)$\\
4 & \multicolumn{2}{l}{$U^TAV$ debe ser diagonal}\\
\hline
\end{tabular}
\end{center}
El paso 3 multiplica el paso 2 por $U^T$ a izquierda y por $V$ a derecha.
Entonces $U^TAV$ conmuta con la matriz diagonal $\Lambda$.
\begin{enumerate}[label=\textup{(\alph*)}]
\item Probar, en el caso $3\times3$, que si $\Sigma\Lambda=\Lambda\Sigma$ con
$\Lambda$ diagonal de elementos \emph{distintos}, entonces $\Sigma$ es
diagonal. Comparar las primeras filas de los dos productos.
\item Identificar dónde falla el argumento si $\Lambda$ tiene autovalores
repetidos, y relacionarlo con la observación de la
Sección~\ref{sec:prueba-SVD} sobre por qué los $\bu$ hay que definirlos como
$A\bv/\sigma$.
\item El mismo problema aparece en las pruebas simples del teorema espectral
$S=Q\Lambda Q^T$. ¿Dónde?
\end{enumerate}

\item\label{ej:parametros} Una matriz de $2\times2$ tiene cuatro elementos y
su SVD, cuatro parámetros: $\theta$, $\phi$, $\sigma_1$, $\sigma_2$. Pasemos
a una de $2\times3$, con seis elementos. ¿Cuántos valores singulares tiene?
La matriz $U$ ($2\times2$) necesita un solo ángulo; ¿cuántos ángulos
determinan entonces a $V$ ($3\times3$)? \emph{Ayuda:} el espacio fila de $A$
es un plano en $\R^3$; hacen falta ciertos ángulos para ubicar ese plano y
uno más para elegir $\bv_1$ y $\bv_2$ dentro de él.

\item Cada matriz de $3\times3$ tiene nueve elementos, así que
$U\Sigma V^T$ debe tener nueve parámetros. ¿Cuántos hay en $U$, en $\Sigma$ y
en $V$? Responder lo mismo para $4\times4$. ¿Cuántos parámetros describen una
rotación del espacio de dimensión $4$?

\item ¿Cuántos números hacen falta para dar la dirección de un vector
unitario $\bv_1$ de $\R^5$? ¿Y la de un unitario $\bv_2$ ortogonal a él? ¿Y
para $\bv_3$, $\bv_4$, $\bv_5$? ¿Cuántos en total?

\item Mostrar que una matriz de $m\times n$ y rango $r$ tiene
$r(m+n-r)$ parámetros libres en su SVD $A=U_r\Sigma_rV_r^T$. ¿Por qué $r$
vectores ortonormales $\bu_1,\dots,\bu_r$ de $\R^m$ tienen
$(m-1)+(m-2)+\cdots+(m-r)$ parámetros? \emph{Otro camino:} usar
$A=CR$ del Capítulo~\ref{cap:algebra-lineal}, donde $R$ contiene una
identidad de $r\times r$ y eso quita $r^2$ parámetros de los $rm+rn$.

\item Si $\bv$ es autovector de $A^TA$ con autovalor $\lambda\neq0$,
¿de qué matriz es autovector $A\bv$, y con qué autovalor?

\item Si $A=U\Sigma V^T$ es cuadrada e invertible, escribir $A^{-1}$ en
términos de $U$, $\Sigma$ y $V$. ¿Cuáles son los valores singulares de
$A^TA$ (no los de $A$)?

\item Si $A=A^T$ tiene columnas ortogonales $\bx_1,\bx_2,\bx_3$ de $\R^3$
con normas $2$, $3$ y $4$, encontrar su SVD.

\item Las razones del éxito de los autovalores están en
$A^k=X\Lambda^kX^{-1}$: (a) los autovalores de $A^k$ son $\lambda_1^k,\dots,
\lambda_n^k$; (b) un autovector de $A$ también lo es de $A^k$. Mostrar que
(a) y (b) son \emph{falsas} para valores y vectores singulares, con
$A=\begin{bsmallmatrix}-2&-6\\ 6&2\end{bsmallmatrix}$.

\end{enumerate}
